\section{Discussion}\label{sec:discussion}

\subsection{What the Findings Mean for Consumer Neuroscience}\label{disc:meaning}

The four findings describe one dataset and one label, but the mechanism behind each is general. The engagement index is a rule we constructed over five gaze statistics and five band-power terms, and the controls show that its gaze terms account for its decodability: a probe on those five statistics recovers the label at ROC-AUC 0.92, a gaze-only LSTM reaches 0.89, and no tested multimodal model, whatever its architecture, exceeded the gaze-only LSTM's ROC-AUC. Because the index is our construction, these are properties of a rule-defined target rather than discoveries about an established measure of engagement; what the audit contributes is the demonstration that the accuracy of a multimodal decoder on such a target can be attributed to the terms the target is built from, and the controls that do so. Any study that reports EEG-based or multimodal decoding of a consumer construct whose label contains gaze statistics, or whose epochs are delimited by gaze as the NeuMa product segments are \cite{ref35}, faces the same question, and the recoverability probe of audit step~1 answers it in seconds. This is the decoding counterpart of a lesson consumer neuroscience learned for its measures: physiological signals earn their place by incremental validity over what simpler measures already capture, tested modality by modality on one footing \cite{ref65}. On this index, no tested EEG encoder showed incremental information over gaze.

The second finding separates two readings that a null on its own cannot. The same recordings, encoders and folds on which no EEG encoder outperformed the EEG-term probe decode resting state versus browsing at 0.85--0.94, so the recordings and pipelines carry usable condition-discriminating information; what has not been found is EEG information about the index beyond the terms it is built from. The control does not isolate neural engagement, since the two conditions differ in visual input, eye movements and muscle activity as well. The single-marker analysis sharpens this: none of the canonical markers, frontal theta, posterior alpha or fronto-posterior phase locking, shows a detectable within-participant difference on the index, and the case studies show frontal theta as a correlate that appears in individual participants and vanishes at the population level. The attribution summaries of the rule-only variant, posterior alpha at O1, O2 and Pz with a horizontal-gaze term in seven of eight rules, are consistent with the attention literature \cite{ref1,ref8} and, more directly, with the five gaze statistics and the P3/C3/F3/Fz band powers from which the label is built. They are learned associations with a rule, and they should be read as such.

The fourth finding places the index against the one behavioural criterion the dataset offers. Whether a product is selected for purchase is associated at ROC-AUC 0.85--0.91 with how long it is looked at over the session, and predicted at 0.48--0.61 by anything inside a 5-s epoch of EEG or gaze; the decoder's 0.61, although significantly above the standard encoders, adds 0.005 to a dwell probe, an increment whose interval includes zero. Selection-related information therefore lives in a behaviour that accumulates across the session, which includes viewing after the selection, rather than in the physiology of a moment inside it; the dataset's own product segments are likewise defined by dwell \cite{ref35}. A protocol that separates looking from choosing, with a prediction cutoff common to all products before the choice, or a self-report collected per stimulus, is the step that would let a physiological decoder be tested prospectively and against engagement rather than against its own definition.

For practice the reading is narrower than the architecture suggested and more useful for being exact. What can be estimated subject-independently on NeuMa is the rule-defined index, and a deployment that has gaze available does not need an EEG pathway to estimate it; where gaze is unavailable, EEG-only accuracy on the index (ROC-AUC 0.59) is not adequate for use. The trained decoder itself is portable, though its preprocessing is offline (Section~\ref{res:robust}): it applies to new participants without per-person calibration, its ranking is threshold-invariant, a label-free per-participant threshold can be recovered online from three unlabelled epochs (Section~\ref{res:robust}), and single-epoch inference runs in 15~ms on a CPU (Section~\ref{sec:system}). Dwell time, available from eye tracking alone, is the selection-related signal these data support.

\subsection{What the Audit Adds}\label{disc:protocol}

Each step of the audit exists because accuracy alone leaves a specific loophole open, and each loophole was live on these data (Table~\ref{tab:loopholes}). An EEG branch that retains a gaze-derived input, compared with baselines trained under a fixed budget, appears significantly better than seven of eight encoders on balanced accuracy (Supplementary Table~S8); gaze-free and against baselines tuned per fold, no significant difference among the nine encoders is detected. A comparison with an untuned eye-tracking encoder shows the decoder leading numerically (Supplementary Table~S8); against the encoder tuned per fold the difference is $-0.011$. An ablation alone credits the dynamic graph with 0.09 of balanced accuracy; density-matched surrogates show that the measured topology was not necessary for that accuracy. Spike sparsity projects an energy advantage; direct timing shows the spiking encoder slower than a dense twin, with no measurable power difference. Rule-activation traces read like explanations; the fidelity test shows a bypass-dominated decision with rule evidence anti-correlated with it. A label whose scaler is fitted on pooled epochs carries an unmeasured dependence on the held-out participant; fold-wise refitting changes nine labels, moves one run, the gaze-only LSTM, by 0.03 in ROC-AUC, and changes no conclusion. None of these checks requires a property of this architecture; all apply to any decoder with the same inputs, and that is the protocol's value. It is stated as eight steps in Section~\ref{sec:audit}, and the code that implements it is released.

\begin{table*}[!htbp]
\caption{Loopholes that evaluation by accuracy alone leaves open, the audit step that closes each, and what the step shows on these data.}\label{tab:loopholes}
\footnotesize
\begin{tabular*}{\textwidth}{@{\extracolsep{\fill}}p{0.32\textwidth}p{0.22\textwidth}p{0.40\textwidth}@{}}
\toprule
What accuracy alone would show & Audit step & What the step shows \\
\midrule
An EEG branch that retains a gaze-derived ROI input beats seven of eight encoders trained under a fixed budget & Gaze-free variant (2); reference encoders on one footing (3) & Gaze-free and against baselines tuned per fold, no significant difference among nine encoders at ROC-AUC 0.49--0.64, none above the EEG-term probe (Section~\ref{res:rq2}) \\
The decoder exceeds an untuned eye-tracking LSTM in ROC-AUC (Supplementary Table~S8) & Reference encoders on one footing (3) & Against the encoder tuned per fold the difference is $-0.011$, $p_{\mathrm{Holm}}=1.0$ (Section~\ref{res:rq1}) \\
Ablating the graph pathway costs 0.09 in balanced accuracy, so the dynamic connectivity appears accuracy-critical & Topology nulls (5) & Density-matched static and random graphs produce no detectable reduction; the measured topology was not necessary for the accuracy (Section~\ref{res:rq3}) \\
Spike sparsity of about 90\% projects an energy advantage & Measured cost (5); positive control (7) & Slower than a dense twin on CPU and GPU with no measurable power difference; 0.55 ROC-AUC alone on a contrast convolutional encoders decode at 0.92--0.94 (Section~\ref{res:rq3}) \\
Rule-activation traces read as explanations of the decision & Rule fidelity (6) & The decision is bypass-dominated ($\alpha=0.57$, agreement 39\%, $r=-0.69$); attribution summaries only from the gate-closed variant, at a cost of 0.04 in ROC-AUC (Sections~\ref{res:rq3} and~\ref{res:rules}) \\
A label scaled and thresholded on pooled epochs carries an unmeasured dependence on the held-out participant & Fold-wise relabelling (4) & Nine of the 385 held-out labels change; no EEG run moves beyond its run-to-run spread and the gaze-only LSTM falls by 0.03 in ROC-AUC, so the dependence is measured and the conclusions stand (Section~\ref{res:rq2}) \\
\botrule
\end{tabular*}
\end{table*}

Three of the steps deserve a word on why they are placed where they are. The linear probe per modality (step~1) is reported before any learned model because it converts label--input coupling into a number: a learned score is read against a probe on the terms it sees, so the coupling becomes a measured quantity rather than a caveat; the probe orders performance and does not identify the information a model uses. The positive control (step~7) is placed beside the null because a null needs a demonstration that the instrument works; it is the difference between ``no usable signal in these recordings'' and ``usable signal, but none found for this index beyond its own terms''. And the topology surrogates (step~5) generalise beyond graphs: any structured component whose input can be replaced by a density-matched surrogate can be tested the same way, which is what an ablation, which removes the processing together with the structure, cannot do, with the same limit in both cases: a retrained surrogate shows what was not necessary, not what the original model used. On reporting, ROC-AUC is used as the headline metric because it is threshold-invariant, because ranking is what a neuromarketing application most often needs, and because per-participant class balance under LOSOCV makes threshold-bound scores volatile; the operating-point metrics are reported beside it. Headline model-comparison scores are means over folds, and pooled values (calibration and probe references) are labelled as such, and the run-to-run spread of the decoder, 0.03 in balanced accuracy and 0.02 in ROC-AUC, is stated beside the decoder's own variants as a reference; it is not a significance threshold, and non-significant differences are reported as undetected rather than as equivalence.

\subsection{Relation to Prior NeuMa Work and Pretrained Encoders}\label{disc:prior}

The only published decoding study on the same recordings combines EEG graph signal processing with gaze dynamics to classify Buy/NoBuy from gaze-defined epochs under within-participant cross-validation and reports Cohen's $\kappa=0.35$ \cite{ref18}. The decoder reaches $\kappa=0.45$ (0.41 at the calibrated operating point) on the engagement index under the harder subject-held-out protocol, and its gaze-free variant $\kappa=0.08$; the targets, epoching and protocols differ, so this is positioning rather than a ranking, and subject-independent ROC-AUC is a quantity prior NeuMa work did not report. Two of the present findings bear on that line of work directly. The benefit of connectivity features that it reports is not reproduced here: with density-matched surrogates the measured connectivity was not necessary for the decoder's accuracy on this index. And because its epochs are gaze-defined and its label behavioural, the dwell-time result of Section~\ref{res:rq4} predicts that a session-level gaze covariate would account for much of any within-epoch decoding of purchase on those data, a prediction the released code can test.

Pretrained EEG transformers trained on thousands of hours of heterogeneous recordings \cite{ref55,ref56}, convolutional attention encoders \cite{ref57} and adaptive tokenisation front-ends \cite{ref58} are the natural next candidates for the EEG pathway, and the audit says where they should be tested. On the index, every encoder tried here, from the linear probe to the dynamic-graph pathway, scored at or below 0.67, and the learning curve plateaus from sixteen participants; a pretrained encoder that exceeded the EEG-term probe would improve on that reference, and showing that it uses information beyond the five terms would further require incremental analyses against flexible term-only models, and one that did not would extend the present result to a much larger pretraining corpus. The positive control, on which plain convolutional encoders already decode at 0.92--0.94, is where a pretrained encoder's advantage would be visible. It is also where attributions of a standard encoder are interpretable: integrated gradients of EEGNet on the resting-state-versus-browsing contrast would show what a reference encoder uses when it succeeds, a complement to the decoder-based explanations of Section~\ref{res:rules}. Substituting such encoders for the EEG pathway within the same protocol is the experiment this study makes answerable.

\subsection{Limitations}\label{disc:limits}

The findings are bounded in the following ways; for each, the reason, what contains it, and what it does not threaten are stated. First, the study rests on one dataset and one label rule that we constructed, because NeuMa is, to our knowledge, the only public resource with synchronised EEG, eye tracking and consumer annotations from a realistic task and it carries no engagement annotation of its own; the predictable properties of the index are those of a rule-defined target, subject-held-out evaluation is the strongest generalisation test available within the dataset, and cross-dataset replication remains the route to external validity. The mechanism of Finding~1, a label that contains gaze statistics, is checkable on any dataset with the recoverability probe. Second, no model is independent of the label, since the EEG encoders can recover its EEG terms and the gaze models its gaze terms; this is why the linear probes are reported and every score is read against them, and why the EEG ranking is anchored on models that receive no gaze. The comparison with a probe orders performance and does not identify which information a model uses, so no information-theoretic limit on the EEG is claimed. Third, the cohort is small for models of this size and each held-out participant contributes about nine epochs, so single-fold metrics are imprecise and cross-validation carries wide error bars \cite{ref63}; the bootstrap intervals, the permutation test, the replicates and the learning curve bound this, and the findings rest on differences of 0.2--0.3 in ROC-AUC or on nulls calibrated by a control at 0.85 or above, not on small margins. Fourth, the comparisons are matched in effort rather than identical: the baselines were tuned per fold over finite, architecture-specific spaces while the decoder's hyperparameters were fixed, and the decoder alone uses a five-member ensemble, a focal loss and augmentation; a larger search could separate the EEG encoders in either direction, and the absence of detected differences in Table~\ref{tab:gaze} is a matched-effort result; the two fusions of standard parts are decision-level only, so the finding that fusing a standard EEG encoder with the gaze model did not exceed it rests on them together with the six attention-based fusion architectures. Fifth, each control was run once in one configuration and, under the pooled-label protocol, the validation participant was drawn independently for the decoder and the baselines, with two folds drawing a single-class participant; the fold-wise runs, which use a matched validation rule that never selects a single-class participant, reproduce the conclusions for the models included in them. The run-to-run spread is a reference for the decoder's own variants, not a significance threshold, and non-significant differences are reported as undetected, not as equivalence, since no equivalence margin was pre-specified and the intervals are wide. The ablations and topology surrogates are retrained models: they show which components were not necessary for the reported accuracy, not which the original trained model used. Sixth, the gaze-free variant was trained without the MMD term as well as without gaze, so its comparison with the decoder removes a regulariser together with the inputs; the MMD ablation ($-0.03$, not significant) is a separate sensitivity check on that confound, not a bound on it. Seventh, the band-power node features rest on 0.5-s windows, so the delta and theta bands are estimated from one to two spectral bins, a coarse resolution for slow bands. Eighth, the product-selection label records intended purchase in a simulated task, not a verified transaction; the epoch anchor is chosen retrospectively over the page view, the dwell covariate accumulates over the whole session including viewing after the selection, so the session-level result is an association rather than a prospective prediction, and the incremental-validity probe stacks held-out probabilities with a dependence that favours the model; a behavioural criterion collected under a protocol that separates looking from choosing, on an independent cohort, is the step that would turn decoding of the index into evidence about engagement. Ninth, the positive control does not isolate neural engagement, and the preprocessing is fitted on whole recordings offline. Tenth, the attribution summaries of the rule-only variant have not been validated as verbal rules, its activations are near-uniform, and rule stability across folds and the sensitivity to the number of rules have not been quantified. None of these limitations changes the direction of any finding: the gaze inputs account for the index's decodability, no tested EEG encoder exceeded the EEG-term probe while the recordings carry condition information, the decoder's proposed mechanisms were not substantiated by performance, and session-level dwell is associated with product selection, under every protocol variant tried.

\subsection{Ethical Considerations}\label{disc:ethics}

Decoding engagement from physiological signals is dual-use. A model that estimates covert engagement could be applied to optimise material for persuasive effect on individuals who have not consented to such profiling, or to infer engagement in settings, such as political messaging, outside the commercial evaluation for which the NeuMa data were collected \cite{ref35}. The present study did not evaluate any such application; the models are trained and assessed on anonymised, consent-collected data for retrospective group-level analysis, and no individual-level, real-time or covert deployment is proposed or endorsed. The findings themselves lower the stakes in one respect, since what the decoder estimates is a gaze-dominated, rule-defined index rather than a covert state, but responsible use of this class of method in applied settings still requires informed participant consent, transparent disclosure of physiological monitoring and independent ethical oversight, consistent with prior calls for ethical governance of consumer neuroscience \cite{ref33}.

\section{Conclusions}\label{sec:conclusion}

This study asked what a multimodal EEG and eye-tracking engagement decoder learns when it is accurate, and answered it on a rule-defined engagement index that we constructed from the NeuMa recordings, by building a representative architecture, comparing it with twenty tuned baselines under subject-held-out evaluation, and auditing it with eight controls anchored on EEGNet, ShallowConvNet and an eye-tracking LSTM. Tuned gaze-only models performed comparably to every tested multimodal system, the decoder included, and fusing EEGNet or ShallowConvNet with the gaze model did not exceed it. No tested EEG encoder outperformed a linear probe on the index's own EEG terms, under the pooled or the fold-wise label, while the same recordings carry condition-discriminating information that convolutional, linear and graph encoders decode at 0.85--0.94. Performance alone did not substantiate the decoder's proposed connectivity, efficiency or interpretability mechanisms: density-matched surrogates showed that the measured topology was not necessary, the spiking front-end was slower than a dense twin, and the rule module's decisions were bypass-dominated, with attribution summaries available from a rule-only variant at a cost of 0.04 in ROC-AUC. On a product-selection label, session-level dwell was associated with selection at 0.85--0.91, and no tested model of a 5-s epoch added detectably to it. These findings do not imply that EEG lacks value for engagement decoding; they show that, on a gaze-containing rule-defined index, the tested EEG encoders and architectural mechanisms did not add to what the index's own terms provide, and they locate the signal that was found. What the study offers beyond the findings is the audit: label recoverability per modality, a gaze-free variant, reference encoders on one footing, a fold-wise label, topology surrogates, a fidelity test, a positive control and a behavioural criterion, each of which closes a loophole that accuracy alone leaves open, and each of which applies to any decoder with the same inputs. Behaviourally grounded labels collected under a protocol that separates looking from choosing, replication on an independent cohort, and pretrained EEG encoders tested within the same protocol are the next steps toward engagement decoding that is more than the recovery of its own definition.
